\section{过去三年的 Language Agent：ReAct、规划、工具使用与开源浪潮}

访谈认为，过去三年 Language Agent 的发展速度超过此前几十年。重要节点包括 ReAct、LLM Planner、Mind2Web、Toolformer、AutoGPT，以及一系列 web agent、computer use agent、robot planning 工作。转写中把 Toolformer 识别成“2-former”，结合语境应指 Meta 的 Toolformer 类工具使用方向；本笔记按访谈语义处理。

\begin{figure}[H]
\centering
\includegraphics[width=0.92\textwidth]{language-agent-stack.png}
\caption{Language Agent 栈：从用户目标到工具、环境反馈与记忆。}
\end{figure}

\begin{knowledgebox}{读图：Language Agent 栈的层次}
最上层是用户目标；往下是自然语言、编程语言和结构化指令；中间是 LLM 的规划与推理；再往下是工具、API、GUI、CLI、代码库等外部接口；底层是环境反馈和记忆。它支持的结论是：Language Agent 不是一个 prompt 技巧，而是一套闭环工程栈。
\end{knowledgebox}

\subsection{ReAct：简单但深远的 insight}

ReAct 的核心是把 reasoning 与 acting 交替起来：模型先推理，再行动，再观察，再继续推理。苏煜强调，很多 Agent 工作看起来技术很简单，但在正确时间点提出正确抽象非常不容易。ReAct 的价值在于，它把 LLM 从单次文本生成推进到环境交互循环。

\noindent\begin{tabularx}{\textwidth}{p{0.22\textwidth}p{0.34\textwidth}X}
\toprule
工作/方向 & 解决的问题 & 与本期主线关系 \\
\midrule
ReAct & 让模型交替进行 reasoning 与 acting & 奠定 LLM 与环境交互的基本模式。 \\
LLM Planner & 用 LLM 做机器人或 embodied planning & 把语言模型用于动作计划，而不只是聊天。 \\
Mind2Web & 构造 web/computer use agent benchmark & 让网页任务成为可评估的 Agent 能力。 \\
Toolformer & 让模型学习何时调用工具 & 连接语言模型与外部 API/工具生态。 \\
AutoGPT & 早期开源 long-horizon agent 项目 & 展示大众对自治代理的想象，也暴露可靠性问题。 \\
\bottomrule
\end{tabularx}

\subsection{从 proof of concept 到资源密集阶段}

苏煜回顾自己从 semantic parsing 转向 language agent 的过程，也解释了为什么越来越多研究者离开学校创业。早期 Agent 工作更像 proof of concept：一个好 idea 可以用低成本方式证明。到 2025 年后，真正有意思的 Agent idea 越来越需要 GPU、API、工程团队和快速试错能力，这与学校资源结构不完全匹配。

\begin{warningbox}{不要把“开源项目火了”误读成“问题解决了”}
AutoGPT 和 OpenClaw 这类项目能迅速聚集注意力，是因为它们展示了可能性。但可能性不等于可靠产品。长期任务中的状态管理、错误恢复、成本控制和安全边界，才是 Agent 真正困难的部分。
\end{warningbox}

\subsection{本章小结}

过去三年的 Language Agent 发展，是从单次文本生成走向环境闭环的过程。ReAct、规划、工具使用、web agent 和开源自治代理共同推动了 OpenClaw Moment 的到来。

